%% TeXora template: Minimal White. Clean white slides, dark text and a single blue accent.
%% Change the colours on the \definecolor lines below; everything else follows them.
% !TEX program = pdflatex
\documentclass[aspectratio=169]{beamer}
\usetheme{default}
\usepackage{amsmath}
\definecolor{ink}{RGB}{31,41,55}
\definecolor{accent}{RGB}{37,99,235}
\definecolor{soft}{RGB}{241,245,249}
\setbeamercolor{normal text}{fg=ink}
\setbeamercolor{structure}{fg=accent}
\setbeamercolor{title}{fg=ink}
\setbeamercolor{subtitle}{fg=accent}
\setbeamercolor{frametitle}{fg=ink}
\setbeamercolor{block title}{bg=soft,fg=accent}
\setbeamercolor{block body}{bg=soft,fg=ink}
\setbeamercolor{block title example}{bg=soft,fg=green!45!black}
\setbeamercolor{block body example}{bg=soft,fg=ink}
\setbeamercolor{block title alerted}{bg=soft,fg=red!70!black}
\setbeamercolor{block body alerted}{bg=soft,fg=ink}
\setbeamerfont{title}{series=\bfseries,size=\LARGE}
\setbeamerfont{frametitle}{series=\bfseries,size=\Large}
\setbeamertemplate{frametitle}{\vspace{0.7em}{\usebeamerfont{frametitle}\usebeamercolor[fg]{frametitle}\insertframetitle}\par\vspace{0.25em}{\color{accent}\rule{3em}{2pt}}\vspace{0.2em}}
\setbeamertemplate{blocks}[rounded]
\setbeamertemplate{itemize items}[circle]
\setbeamertemplate{navigation symbols}{}
\setbeamertemplate{footline}{\hspace{1.4em}{\color{ink!45}\footnotesize\insertshorttitle}\hfill{\color{ink!45}\footnotesize\insertframenumber}\hspace{1.4em}\vspace{0.9em}}
\setbeameroption{show notes on second screen=right}

\title{Your Presentation Title}
\subtitle{A short subtitle}
\author{Your Name}
\institute{Your Institution}
\date{\today}

\begin{document}

\begin{frame}
  \titlepage
  \note{Welcome everyone. Say who you are and what this talk covers.}
\end{frame}

\begin{frame}{Outline}
  \tableofcontents
  \note{Give the plan in one sentence.}
\end{frame}

\section{Introduction}
\begin{frame}{Key Ideas}
  \begin{itemize}
    \item<1-> Your first point
    \item<2-> Your second point, shown on the next step
    \item<3-> Your third point
  \end{itemize}
  \vspace{0.6em}
  \begin{block}<3->{Takeaway}
    One sentence the audience should remember.
  \end{block}
  \note{Each point appears on its own step. Pause after the takeaway.}
\end{frame}

\section{Details}
\begin{frame}{Two Columns}
  \begin{columns}[T]
    \begin{column}{0.47\textwidth}
      \begin{exampleblock}{Example}
        A short example, a result or a figure.
      \end{exampleblock}
    \end{column}
    \begin{column}{0.47\textwidth}
      \begin{alertblock}{Watch out}
        A caveat or a common mistake.
      \end{alertblock}
    \end{column}
  \end{columns}
  \note{Compare the two sides, left to right.}
\end{frame}

\begin{frame}{An Equation}
  \begin{equation*}
    f(x) = \sum_{k=0}^{n} a_k\, x^k
  \end{equation*}
  Explain what each symbol means, and why it matters here.
  \note{Read the equation aloud once, slowly.}
\end{frame}

\section{Conclusion}
\begin{frame}{Summary}
  \begin{enumerate}
    \item Your first conclusion
    \item Your second conclusion
  \end{enumerate}
  \vfill
  \begin{center}
    {\Large Thank you. Questions?}
  \end{center}
  \note{Thank the audience and invite questions.}
\end{frame}

\end{document}
